\section{\HMTLang{Composición de acción, frecuencia propia y reciprocidad radial}{Composition of action, proper frequency, and radial reciprocity}}
\label{delta22:vi:composicion}
\HMTLang{El operador de masa precedente recibe la firma, el carácter y las torres de la historia enriquecida. La sección de acción fija su escala absoluta y participa también en el contraángulo y en el corrector de retorno. La composición siguiente hace explícita esa concurrencia: una misma coordenada estructural de ruta determina la masa, la frecuencia propia, dos lecturas radiales recíprocas y, en la realización canónica especificada, la ponderación térmica.}{The preceding mass operator receives the signature, character, and towers of the enriched history. The action section fixes its absolute scale and also participates in the counterangle and return correction. The following composition makes that concurrence explicit: one structural route coordinate determines mass, proper frequency, two reciprocal radial readings, and, in the specified canonical realization, thermal weighting.}
\HMTLang{Se conserva el orden APP--TRIT--TPK, estado enriquecido y estructura discreta conjunta del continuo; las coordenadas regionales y \(\alpha\) son salidas previamente generadas. La selección de \(K\), la incidencia y la orientación permanecen en sus antecedentes. Las operaciones de este apartado son composiciones posteriores de los lectores ya construidos.}{The order APP--TRIT--TPK, enriched state, and joint discrete structure of the continuum is retained; the regional coordinates and \(\alpha\) are previously generated outputs. The selection of \(K\), incidence, and orientation remain in their antecedents. The operations in this section are subsequent compositions of the readers already constructed.}

\subsection{\HMTLang{Restitución radial de la sección de acción}{Radial recovery of the action section}}
\HMTLang{Con el par angular en una misma unidad, escribimos}{With the angular pair in the same unit, write}
\[
 A=1000\alpha,\quad C^*=2(\eta_{\rm ret}+\alpha),\quad
 \xi=C^*/A,\quad
 R_\star=\left(\frac{1-\xi}{1+\xi}\right)^{1/4},\quad q=R_\star^4.
\]
\HMTLang{En la rama positiva \(0<\eta_{\rm ret}<499\alpha\), se tiene \(0<q<499/501\). La inversión explícita es}{On the positive branch \(0<\eta_{\rm ret}<499\alpha\), one has \(0<q<499/501\). The explicit inverse is}
\begin{equation}
 \eta_{\rm ret}=\alpha\frac{499-501q}{1+q},\qquad
 \hbar_{\rm ret}=10^{-34}\mathcal U_S\eta_{\rm ret},\qquad
 R_{\rm act}=\frac{H_5(1+q)}{\alpha(499-501q)}.
 \label{delta22:vi:accion}
\end{equation}
\begin{proof}
\HMTLang{La identidad \(q=(1-\xi)/(1+\xi)\) da \(\xi=(1-q)/(1+q)\). Sustituirla en \(\eta_{\rm ret}=\alpha(500\xi-1)\) prueba la primera fórmula; las otras dos usan las definiciones de la sección de acción y del corrector \(R_{\rm act}=H_5/\eta_{\rm ret}\). Para las dos orientaciones, sea \(\varepsilon=\pm1\) y conserve \(q_\varepsilon=(1-\varepsilon\xi)/(1+\varepsilon\xi)\). Entonces}{The identity \(q=(1-\xi)/(1+\xi)\) gives \(\xi=(1-q)/(1+q)\). Substitution into \(\eta_{\rm ret}=\alpha(500\xi-1)\) proves the first formula; the other two use the definitions of the action section and correction \(R_{\rm act}=H_5/\eta_{\rm ret}\). For both orientations, let \(\varepsilon=\pm1\) and retain \(q_\varepsilon=(1-\varepsilon\xi)/(1+\varepsilon\xi)\). Then}
\[
 \eta_{\rm ret}=\alpha\left[500\varepsilon
 \frac{1-q_\varepsilon}{1+q_\varepsilon}-1\right].
\]
\HMTLang{La transformación \((\varepsilon,q_\varepsilon)\mapsto(-\varepsilon,q_\varepsilon^{-1})\) conserva la expresión y, con ella, \(\hbar_{\rm ret}\) y \(R_{\rm act}\).}{The transformation \((\varepsilon,q_\varepsilon)\mapsto(-\varepsilon,q_\varepsilon^{-1})\) preserves the expression and hence \(\hbar_{\rm ret}\) and \(R_{\rm act}\).}
\end{proof}
\HMTLang{Esta restitución conserva la marca de orientación, \(\alpha\), \(H_5\) y la base dimensional. La firma y las torres continúan procediendo del registro de ruta.}{This recovery retains the orientation mark, \(\alpha\), \(H_5\), and the dimensional basis. The signature and towers continue to arise from the route record.}

\subsection{\HMTLang{Sección longitudinal, reloj inducido y cotransformación electrónica}{Longitudinal section, induced clock, and electronic cotransformation}}
\label{delta22:vi:longitud}
\HMTLang{Sean \(L_\star>0\) la sección longitudinal de referencia y \(c_{\rm int}>0\) la velocidad constitutiva en la coordenatización elegida. Escribimos \(\alpha_{\rm HMT}\) y \(\pi_{\rm HMT}\) para las coordenadas ya generadas por el estado enriquecido; no son datos metrológicos que seleccionen esta construcción. En las fórmulas previamente asentadas que siguen, \(\alpha:=\alpha_{\rm HMT}\) y \(\pi:=\pi_{\rm HMT}\). La incidencia marcada contiene}{Let \(L_\star>0\) be the reference longitudinal section and \(c_{\rm int}>0\) the constitutive velocity in the chosen chart. We write \(\alpha_{\rm HMT}\) and \(\pi_{\rm HMT}\) for the coordinates already generated by the enriched state; they are not metrological data selecting this construction. In the previously established formulae below, \(\alpha:=\alpha_{\rm HMT}\) and \(\pi:=\pi_{\rm HMT}\). The marked incidence contains}
\begin{equation}
 N=12\cdot4\cdot120=5760,
 \qquad r_N=\frac{N-1}{4N}=\frac{5759}{23040},
 \qquad
 L_\alpha=\alpha_{\rm HMT}^{16}r_NL_\star .
 \label{delta22:vi:Lalpha}
\end{equation}
\HMTLang{El coeficiente \(r_N\) es la expectativa diagonal del retorno residual de Hadamard, \(\Delta(CC^*)=r_NI\). El canal complementario \(1/(4N)\) se conserva como complemento, pero no se suma al observable longitudinal. Por tanto, \(L_\alpha\) queda construido antes del reloj y antes de cualquier lectura gravitatoria.}{The coefficient \(r_N\) is the diagonal expectation of the residual Hadamard return, \(\Delta(CC^*)=r_NI\). The complementary channel \(1/(4N)\) is retained as a complement but is not added to the longitudinal observable. Hence \(L_\alpha\) is constructed before the clock and before any gravitational reading.}

\HMTLang{El reloj de \(108\) pasos es la realización circular posterior de esa longitud:}{The \(108\)-step clock is the subsequent circular realization of that length:}
\begin{equation}
 t_0=\frac{\pi_{\rm HMT}L_\alpha}{54c_{\rm int}},\qquad
 \omega_0=\frac{2\pi_{\rm HMT}}{108t_0}=\frac{c_{\rm int}}{L_\alpha},
 \qquad R_{108}=\frac{c_{\rm int}}{\omega_0}=L_\alpha .
 \label{delta22:vi:reloj-longitudinal}
\end{equation}
\HMTLang{Fijada una sección de acción retornada \(\hbar_{\rm ret}>0\), se definen la energía y la masa del reloj, y sólo entonces el factor electrónico:}{For a fixed returned action section \(\hbar_{\rm ret}>0\), the clock energy and mass are defined first, and only then the electronic factor:}
\begin{equation}
 Q_L=\frac{\hbar_{\rm ret}c_{\rm int}}{L_\alpha},\qquad
 E_{\rm clk}=Q_L,\qquad
 m_{\rm clk}=\frac{Q_L}{c_{\rm int}^2},\qquad
 E_e^{(0)}=\varepsilon_e^{(0)}U_E,\qquad
 b_e(Q_L)=\frac{E_e^{(0)}}{Q_L}.
 \label{delta22:vi:escala-longitudinal}
\end{equation}
\HMTLang{En lo que sigue, \(b_e\) abrevia la función evaluada \(b_e(Q_L)\).}{Hereafter, \(b_e\) abbreviates the evaluated function \(b_e(Q_L)\).}
\HMTLang{Se conserva uno de los lectores de retorno \(r\in\{D,\Omega\}\). El carácter refinado es el construido en el apartado precedente, con convergencia y torres relativas al electrón conservadas.}{One of the return readers \(r\in\{D,\Omega\}\) is retained. The refined character is the one constructed in the preceding section, retaining convergence and towers relative to the electron.}
\HMTLang{La coordenada electrónica basal y su exponente, derivados anteriormente, son}{The basal electronic coordinate and its exponent, derived above, are}
\[
 d_4=\frac{\pi-e\log\pi}{270},\quad
 \Delta_4=d_4\left(1+\frac\pi{729}\right),\quad
 \varepsilon_e^{(0)}=\frac{\sqrt3}{4}
 \left[e^{\varphi/\pi^2}-22\alpha^3\right](1+15\Delta_4),\quad
 \kappa_e=80-54\alpha+6\Delta_4.
\]
\HMTLang{Para una ruta con realización escalar positiva,}{For a route with a positive scalar realization,}
\begin{equation}
 \Xi_\Gamma=b_e\,
 \chi_{\rm ref}(\sigma_\Gamma-\sigma_e,\eta_\Gamma-\eta_e)
 R_{\rm act}^{\kappa_\Gamma^r},\qquad
 m_\Gamma^\beta=m_{\rm clk}\Xi_\Gamma.
 \label{delta22:vi:xi}
\end{equation}
\HMTLang{En esta fórmula se emplean las torres absolutas antes de tomar su diferencia; cuando el registro ya proporciona las torres relativas, esa diferencia se toma una sola vez. El término de altura \(120\) conserva asimismo su única inclusión en el carácter. En el estado central,}{This formula uses absolute towers before taking their difference; when the record already provides relative towers, that difference is taken only once. The height-\(120\) term likewise retains its single inclusion in the character. At the central state,}
\[
 m_e^\beta=\frac{\varepsilon_e^{(0)}U_E}{c_{\rm int}^2}
 R_{\rm act}^{\kappa_e}.
\]
\HMTLang{La igualdad se comprueba cancelando \(Q_L\) entre \(m_{\rm clk}\) y \(b_e(Q_L)\). Así se conserva la normalización electrónica completa. La acción interviene en la escala absoluta, en \(R_{\rm act}\) y en el contraángulo que determina los canales constitutivos. No aparece ninguna corrección decimal dependiente de la especie: la década de la sección de acción permanece exclusivamente en la rama absoluta acción--reloj--masa.}{The equality follows by cancelling \(Q_L\) between \(m_{\rm clk}\) and \(b_e(Q_L)\). This retains the complete electronic normalization. Action enters the absolute scale, \(R_{\rm act}\), and the counterangle determining the constitutive channels. No species-dependent decimal correction occurs: the decade of the action section remains exclusively in the absolute action--clock--mass branch.}

\HMTLang{La covariancia de coordenatización impide congelar \(b_e\). Con \(c_{\rm int}\), \(E_e^{(0)}\) y todos los datos adimensionales de la ruta fijados ---carácter refinado, exponentes y \(R_{\rm act}\)---, la cotransformación de las secciones longitudinal y de acción se escribe}{Chart covariance prevents \(b_e\) from being frozen. With \(c_{\rm int}\), \(E_e^{(0)}\), and all dimensionless route data fixed---the refined character, exponents, and \(R_{\rm act}\)---, the cotransformation of the longitudinal and action sections is written}
\[
 L_\star\longmapsto\lambda L_\star,\qquad
 \hbar_{\rm ret}\longmapsto\mu\hbar_{\rm ret},\qquad \lambda,\mu>0,
\]
\HMTLang{entonces}{then}
\begin{equation}
 L_\alpha\longmapsto\lambda L_\alpha,\qquad
 Q_L,m_{\rm clk}\longmapsto\frac{\mu}{\lambda}(Q_L,m_{\rm clk}),
 \qquad b_e\longmapsto\frac{\lambda}{\mu}b_e,\qquad
 m_{\rm clk}b_e=\frac{E_e^{(0)}}{c_{\rm int}^2}.
 \label{delta22:vi:be-cotransformacion}
\end{equation}
\HMTLang{Por consiguiente, la variación de secciones cotransforma el factor electrónico y no altera las masas ya construidas ni sus tablas.}{Consequently, varying the sections cotransforms the electronic factor and does not alter the already constructed masses or their tables.}

\HMTLang{En las identidades siguientes se abrevia \(m_\Gamma=m_\Gamma^\beta\), conservando la normalización electrónica completa.}{In the following identities, \(m_\Gamma=m_\Gamma^\beta\) is used as an abbreviation, retaining the complete electronic normalization.}
\begin{proposition}[\HMTLang{Frecuencia propia y longitud reducida}{Proper frequency and reduced wavelength}]
\HMTLang{Con \(E_\Gamma=m_\Gamma c_{\rm int}^2\), \(\Omega_\Gamma=E_\Gamma/\hbar_{\rm ret}\), \(\nu_\Gamma=\Omega_\Gamma/(2\pi_{\rm HMT})\) y \(R_{108}=L_\alpha\), se tiene}{With \(E_\Gamma=m_\Gamma c_{\rm int}^2\), \(\Omega_\Gamma=E_\Gamma/\hbar_{\rm ret}\), \(\nu_\Gamma=\Omega_\Gamma/(2\pi_{\rm HMT})\), and \(R_{108}=L_\alpha\), one has}
\[
 E_\Gamma=E_{\rm clk}\Xi_\Gamma,\quad
 \Omega_\Gamma=\omega_0\Xi_\Gamma,\quad
 \nu_\Gamma=\frac{\Xi_\Gamma}{108t_0},\quad
 \bar\lambda_\Gamma=\frac{\hbar_{\rm ret}}{m_\Gamma c_{\rm int}}
 =\frac{R_{108}}{\Xi_\Gamma}.
\]
\end{proposition}
\begin{proof}
\HMTLang{Sustituir \eqref{delta22:vi:xi} y las definiciones de \(m_{\rm clk}\) y \(E_{\rm clk}\) produce las cuatro identidades. Para dos rutas positivas, el cociente de frecuencias es el cociente de masas; la propiedad multiplicativa del carácter conserva la firma y las torres relativas, junto con \(R_{\rm act}^{\kappa_Y^r-\kappa_X^r}\).}{Substituting \eqref{delta22:vi:xi} and the definitions of \(m_{\rm clk}\) and \(E_{\rm clk}\) gives all four identities. For two positive routes, the frequency ratio equals the mass ratio; the multiplicative property of the character retains the relative signature and towers together with \(R_{\rm act}^{\kappa_Y^r-\kappa_X^r}\).}
\end{proof}

\subsection{\HMTLang{Rigidez conjunta de acción y lectura radial}{Joint rigidity of action and radial reading}}
\HMTLang{En esta subsección abreviamos \(\hbar=\hbar_{\rm ret}\). Sea \(M=M^{\beta,r}\) el operador de masa ya construido sobre una fibra masiva finita y estrictamente positiva. Definimos, después de \(M\),}{In this subsection we abbreviate \(\hbar=\hbar_{\rm ret}\). Let \(M=M^{\beta,r}\) be the already constructed mass operator on a finite, strictly positive massive fibre. After \(M\), define}
\[
 X=\frac{M}{m_{\rm clk}},\qquad
 H=Mc_{\rm int}^2=Q_LX,\qquad
 \Lambda=L_\alpha X^{-1},\qquad
 \mathcal R_g=L_\alpha X.
\]
\begin{theorem}[\HMTLang{Dos relaciones conservadas y coeficiente gravitatorio único}{Two conserved relations and unique gravitational coefficient}]
\HMTLang{Las lecturas anteriores satisfacen}{The preceding readings satisfy}
\begingroup
\renewcommand{\theHequation}{delta22.R01}
\begin{equation}
 H\Lambda=\hbar c_{\rm int}I.
 \tag{R01}\label{delta22:vi:R01}
\end{equation}
\endgroup
\begingroup
\renewcommand{\theHequation}{delta22.R02}
\begin{equation}
 \mathcal R_g\Lambda=\Lambda\mathcal R_g=L_\alpha^2I.
 \tag{R02}\label{delta22:vi:R02}
\end{equation}
\endgroup
\HMTLang{Si la realización física identifica posteriormente la lectura radial reducida mediante \(\mathcal R_g=(G/c_{\rm int}^2)M\), entonces existe un único coeficiente compatible:}{If the physical realization subsequently identifies the reduced radial reading by \(\mathcal R_g=(G/c_{\rm int}^2)M\), then there is a unique compatible coefficient:}
\begingroup
\renewcommand{\theHequation}{delta22.G.unica}
\begin{equation}
 \boxed{G=\frac{c_{\rm int}^3L_\alpha^2}{\hbar}.}
 \label{delta22:vi:G-unica}
\end{equation}
\endgroup
\end{theorem}
\begin{proof}
\HMTLang{Las definiciones y \(Q_L=\hbar c_{\rm int}/L_\alpha\) dan directamente (R01) y (R02). De ellas se obtiene \(\mathcal R_g=L_\alpha^2H/(\hbar c_{\rm int})\). Sustituir \(H=Mc_{\rm int}^2\) y compararlo con \(\mathcal R_g=(G/c_{\rm int}^2)M\) produce \eqref{delta22:vi:G-unica}. Si \(G'\) satisface las mismas lecturas, \((G'-G)M/c_{\rm int}^2=0\); la invertibilidad de \(M\) fuerza \(G'=G\). El sector nulo se trata por separado y no se invierte. En particular, \(G\) se deduce después del operador de masa y no interviene en su construcción.}{The definitions and \(Q_L=\hbar c_{\rm int}/L_\alpha\) give (R01) and (R02) directly. They imply \(\mathcal R_g=L_\alpha^2H/(\hbar c_{\rm int})\). Substitution of \(H=Mc_{\rm int}^2\) and comparison with \(\mathcal R_g=(G/c_{\rm int}^2)M\) yield \eqref{delta22:vi:G-unica}. If \(G'\) satisfies the same readings, then \((G'-G)M/c_{\rm int}^2=0\); invertibility of \(M\) forces \(G'=G\). The null sector is treated separately and is not inverted. In particular, \(G\) is deduced after the mass operator and does not enter its construction.}
\end{proof}

\HMTLang{Si \(X\) es sólo semidefinido, se restringe a \((\ker X)^\perp\). Con la pseudoinversa \(X^+\), \(\Lambda^+=L_\alpha X^+\) y \(P=P_{(\ker X)^\perp}\), las relaciones se convierten en \(H\Lambda^+=\hbar c_{\rm int}P\) y \(\mathcal R_g\Lambda^+=L_\alpha^2P\). El sector nulo conserva su tratamiento propio y no produce una determinación vacua de \(G\).}{If \(X\) is only positive semidefinite, the construction is restricted to \((\ker X)^\perp\). With the pseudoinverse \(X^+\), \(\Lambda^+=L_\alpha X^+\), and \(P=P_{(\ker X)^\perp}\), the relations become \(H\Lambda^+=\hbar c_{\rm int}P\) and \(\mathcal R_g\Lambda^+=L_\alpha^2P\). The null sector retains its separate treatment and does not yield a vacuous determination of \(G\).}

\HMTLang{Para cualquier sección de acción positiva \(\hbar_a\), sea \(Q_{L,a}:=\hbar_a c_{\rm int}/L_\alpha=E_{P,a}\) la energía longitudinal de Planck en esa misma sección. La misma prueba da \(G_a=c_{\rm int}^3L_\alpha^2/\hbar_a\) y conserva \(\hbar_aG_a/c_{\rm int}^3=L_\alpha^2\). La representación mediante reloj es un corolario, no la premisa de la deducción:}{For any positive action section \(\hbar_a\), let \(Q_{L,a}:=\hbar_a c_{\rm int}/L_\alpha=E_{P,a}\) be the longitudinal Planck energy in that same section. The same proof gives \(G_a=c_{\rm int}^3L_\alpha^2/\hbar_a\) and preserves \(\hbar_aG_a/c_{\rm int}^3=L_\alpha^2\). The clock representation is a corollary, not the premise of the deduction:}
\begin{corollary}[\HMTLang{Forma circular}{Circular form}]
\begin{equation}
 t_0=\frac{\pi_{\rm HMT}L_\alpha}{54c_{\rm int}}
 \quad\Longrightarrow\quad
 G_a=\left(\frac{54}{\pi_{\rm HMT}}\right)^2
 \frac{c_{\rm int}^5t_0^2}{\hbar_a}.
 \label{delta22:vi:G-circular}
\end{equation}
\end{corollary}

\begin{theorem}[\HMTLang{Reciprocidad de las lecturas radiales}{Reciprocity of radial readings}]
\HMTLang{Con la convención reducida \(r_{g,\Gamma}^{(a)}=G_am_\Gamma/c_{\rm int}^2\),}{With the reduced convention \(r_{g,\Gamma}^{(a)}=G_am_\Gamma/c_{\rm int}^2\),}
\[
 r_{g,\Gamma}^{(a)}=\frac{\hbar_{\rm ret}}{\hbar_a}R_{108}\Xi_\Gamma,
 \qquad
 r_{g,\Gamma}^{(a)}\bar\lambda_\Gamma
 =\frac{\hbar_{\rm ret}}{\hbar_a}R_{108}^2.
\]
\HMTLang{En la misma sección \(a={\rm ret}\),}{In the same section \(a={\rm ret}\),}
\begin{equation}
 r_{g,\Gamma}=R_{108}\Xi_\Gamma,\qquad
 \bar\lambda_\Gamma=R_{108}/\Xi_\Gamma,\qquad
 r_{g,\Gamma}\bar\lambda_\Gamma=R_{108}^2.
 \label{delta22:vi:radios}
\end{equation}
\end{theorem}
\begin{proof}
\HMTLang{Sustituir \(G_a\) y \(m_\Gamma=\hbar_{\rm ret}\omega_0\Xi_\Gamma/c_{\rm int}^2\) da \(r_g^{(a)}=(\hbar_{\rm ret}/\hbar_a)R_{108}^2\omega_0\Xi_\Gamma/c_{\rm int}\). La identidad \(R_{108}\omega_0=c_{\rm int}\) prueba la primera fórmula; multiplicar por \(R_{108}/\Xi_\Gamma\) da la segunda. La elección de una misma sección cancela el cociente de acciones.}{Substituting \(G_a\) and \(m_\Gamma=\hbar_{\rm ret}\omega_0\Xi_\Gamma/c_{\rm int}^2\) gives \(r_g^{(a)}=(\hbar_{\rm ret}/\hbar_a)R_{108}^2\omega_0\Xi_\Gamma/c_{\rm int}\). The identity \(R_{108}\omega_0=c_{\rm int}\) proves the first formula; multiplication by \(R_{108}/\Xi_\Gamma\) gives the second. Choosing the same section cancels the action ratio.}
\end{proof}
\HMTLang{La convención \(2Gm/c^2\) introduciría un factor dos en el producto y constituiría otro lector. Aquí las longitudes son lecturas del mismo estado: en coordenadas logarítmicas sus desplazamientos son \(\log\Xi_\Gamma\) y \(-\log\Xi_\Gamma\). Su media geométrica es \(R_{108}=L_\alpha\). La identificación física de \(\mathcal R_g\) se declara expresamente en el teorema anterior; no redefine el operador másico.}{The convention \(2Gm/c^2\) would introduce a factor of two into the product and would constitute a different reader. Here the lengths are readings of the same state: in logarithmic coordinates their displacements are \(\log\Xi_\Gamma\) and \(-\log\Xi_\Gamma\). Their geometric mean is \(R_{108}=L_\alpha\). The physical identification of \(\mathcal R_g\) is stated explicitly in the preceding theorem; it does not redefine the mass operator.}

\subsection{\HMTLang{Lectura térmica del espectro y compatibilidad conjunta}{Thermal reading of the spectrum and joint compatibility}}
\HMTLang{La transducción térmica es una aplicación lineal positiva entre las rectas dimensionales orientadas de temperatura y energía: \(k_B:\mathcal L_\Theta\to\mathcal L_E\). La distribución uniforme de los tres estados locales tiene información \(-3(1/3)\log(1/3)=\log3\). Una sección positiva de temperatura \(\Theta_{108}\), normalizada por la energía de una decisión cronológica, determina de manera única el transductor mediante}{Thermal transduction is a positive linear map between the oriented dimensional lines of temperature and energy: \(k_B:\mathcal L_\Theta\to\mathcal L_E\). The uniform distribution over the three local states has information \(-3(1/3)\log(1/3)=\log3\). A positive temperature section \(\Theta_{108}\), normalized by the energy of one chronological decision, uniquely determines the transducer through}
\[
 k_B\Theta_{108}\log3=E_{\rm clk}.
\]
\HMTLang{En bases positivas \(u_E,u_\Theta\), sean \(E_{\rm clk}=\epsilon u_E\), \(\Theta_{108}=\theta u_\Theta\) y \(k_B(u_\Theta)=b u_E\). La ecuación da \(b=\epsilon/(\theta\log3)>0\), que prueba existencia y unicidad. Bajo \(u'_E=s_Eu_E\), \(u'_\Theta=s_\Theta u_\Theta\), se tiene \(b'=(s_\Theta/s_E)b\), \(\epsilon'=\epsilon/s_E\), \(\theta'=\theta/s_\Theta\); la identidad se conserva. La coordenada individual del transductor mantiene esas bases y la sección térmica como datos declarados. La acción y el período fijan previamente \(E_{\rm clk}=h_{\rm ret}/(108t_0)\).}{In positive bases \(u_E,u_\Theta\), let \(E_{\rm clk}=\epsilon u_E\), \(\Theta_{108}=\theta u_\Theta\), and \(k_B(u_\Theta)=b u_E\). The equation gives \(b=\epsilon/(\theta\log3)>0\), proving existence and uniqueness. Under \(u'_E=s_Eu_E\), \(u'_\Theta=s_\Theta u_\Theta\), one has \(b'=(s_\Theta/s_E)b\), \(\epsilon'=\epsilon/s_E\), and \(\theta'=\theta/s_\Theta\); the identity is preserved. The individual coordinate of the transducer retains those bases and the thermal section as declared data. Action and period previously fix \(E_{\rm clk}=h_{\rm ret}/(108t_0)\).}
\HMTLang{En la realización canónica a esa temperatura, con origen energético común y potencial químico nulo,}{In the canonical realization at that temperature, with a common energy origin and zero chemical potential,}
\[
 \beta_{108}=(k_B\Theta_{108})^{-1},\qquad
 \beta_{108}E_\Gamma=(\log3)\Xi_\Gamma,\qquad
 e^{-\beta_{108}E_\Gamma}=3^{-\Xi_\Gamma}.
\]
\HMTLang{La prueba sustituye \(E_\Gamma=E_{\rm clk}\Xi_\Gamma\). Para la fibra finita, su expresión operatoria completa es}{The proof substitutes \(E_\Gamma=E_{\rm clk}\Xi_\Gamma\). For the finite fibre, its full operator expression is}
\[
 S(\rho):=-\operatorname{Tr}(\rho\log\rho),\qquad
 Z=\operatorname{Tr}(3^{-X}),\qquad
 \rho_{108}=3^{-X}/Z.
\]
\[
 S(\rho_{108})=\log Z+(\log3)\operatorname{Tr}(\rho_{108}X).
\]
\HMTLang{La última igualdad procede de \(\log\rho_{108}=-(\log3)X-(\log Z)I\) y \(\operatorname{Tr}\rho_{108}=1\). Si los autovalores se agrupan con multiplicidades físicas \(g_\Gamma\), \(Z=\sum_\Gamma g_\Gamma3^{-\Xi_\Gamma}\). Se cuentan estados de la realización; las descripciones equivalentes de una historia conservan su identificación. En dimensión infinita se exige \(3^{-X}\) de clase traza y, para entropía finita, \(\operatorname{Tr}(X3^{-X})<\infty\).}{The last equality follows from \(\log\rho_{108}=-(\log3)X-(\log Z)I\) and \(\operatorname{Tr}\rho_{108}=1\). If eigenvalues are grouped with physical multiplicities \(g_\Gamma\), then \(Z=\sum_\Gamma g_\Gamma3^{-\Xi_\Gamma}\). States of the realization are counted; equivalent descriptions of a history retain their identification. In infinite dimension, \(3^{-X}\) must be trace class and finite entropy additionally requires \(\operatorname{Tr}(X3^{-X})<\infty\).}

\HMTLang{Aquí \(S\) es la entropía adimensional de von Neumann; la entropía dimensional se obtiene mediante el transductor \(k_B\).}{Here \(S\) is dimensionless von Neumann entropy; dimensional entropy is obtained through the transducer \(k_B\).}

\HMTLang{La composición completa, en una misma sección de acción y usando \(R_{108}=L_\alpha\), queda}{The full composition, on the same action section and using \(R_{108}=L_\alpha\), is}
\begin{equation}
 \boxed{\Xi_\Gamma=\frac{m_\Gamma}{m_{\rm clk}}
 =\frac{\Omega_\Gamma}{\omega_0}
 =\frac{r_{g,\Gamma}}{L_\alpha}
 =\frac{L_\alpha}{\bar\lambda_\Gamma}
 =\frac{E_\Gamma}{k_B\Theta_{108}\log3}.}
 \label{delta22:vi:compatibilidad}
\end{equation}
\HMTLang{Esta igualdad identifica las lecturas de una misma sección; la genealogía que la produce permanece en \(b_e\), el carácter refinado y el exponente de retorno. En \(\Xi=1\), las lecturas radiales coinciden y la energía es \(E_{\rm clk}\); la existencia de una especie en ese punto depende del espectro efectivamente construido. La identidad conjunta proporciona un control de compatibilidad entre las realizaciones dimensionales y conserva la distinción entre energía propia, acción, temperatura y memoria de ruta.}{This equality identifies readings of the same section; the genealogy producing it remains in \(b_e\), the refined character, and the return exponent. At \(\Xi=1\), the radial readings coincide and the energy is \(E_{\rm clk}\); whether a species occupies that point depends on the spectrum actually constructed. The joint identity provides a compatibility check between dimensional realizations and retains the distinction between proper energy, action, temperature, and route memory.}
